% Source: https://artofproblemsolving.com/community/c3h2425117_100th_post
% Stable identity: aops-topic-2425117-q02
At an intersection point both equations hold, so
\[x^2+bx+c=x^2+cx+b\iff(b-c)x=b-c.\]
\begin{itemize}
\item If $b\ne c$, then $x=1$ and $y=1+b+c$. The graphs meet in exactly one point, $\boxed{(1,\,1+b+c)}$.
\item If $b=c$, the two equations are the same, so the graphs coincide and every point of the parabola $y=x^2+bx+b$ is an intersection point.
\end{itemize}
For example, with $b=1$ and $c=-1$ the parabolas $y_1=x^2+x-1$ and $y_2=x^2-x+1$ meet only at $(1,1)$, as in the sketch.

\begin{asy}
usepackage("amsmath");
size(5cm);
defaultpen(fontsize(10pt));
import graph;
real f(real x){return x^2+x-1;} real g(real x){return x^2-x+1;}
draw((-2.6,0)--(2.6,0),Arrow(6)); draw((0,-1.6)--(0,6.5),Arrow(6));
label("$x$",(2.6,0),E); label("$y$",(0,6.5),N);
draw(graph(f,-2.4,2.05),blue+linewidth(0.9)); draw(graph(g,-1.95,2.4),red+linewidth(0.9));
dot((1,1)); label("$(1,1)$",(1,1),SE);
label("$y_1$",(2.05,f(2.05)),N,blue); label("$y_2$",(-1.95,g(-1.95)),N,red);
\end{asy}
